\documentclass[11pt,a4paper]{article}

\usepackage[utf8]{inputenc}
\usepackage[T1]{fontenc}
\usepackage{lmodern}
\usepackage{geometry}
\usepackage{hyperref}
\usepackage{setspace}

\geometry{margin=2.5cm}
\setstretch{1.08}
\hypersetup{hidelinks}

\title{Rapport intermediaire -- \v{E}tat d'avancement}
\author{}
\date{}

\begin{document}
\maketitle

\section{\'Etat actuel du projet}

\subsection{Cr\'eation, ex\'ecution et terminaison}


\subsection{Synchronisation}

\subsection{preemption}

\section{Pourquoi certains tests sont meilleurs que pthreads}


\subsection{Moins de travail c\^ot\'e noyau}


\subsection{Ordonnancement adapt\'e au cas d'usage}


\subsection{Co\^ut de cr\'eation ma\^itris\'e}


\section{Pourquoi pthreads restent meilleurs sur certains tests}


\subsection{Les pthreads b\'en\'eficient de l'optimisation syst\`eme}


\subsection{Les calculs longs favorisent le parall\'elisme r\'eel}


\subsection{Certaines charges sont trop petites}


\section{Tests encore limit\'es ou non prioritaires \`a ce stade}

\section{Ce qui sera fait ensuite}

La suite du travail consiste \`a aller vers les objectifs avanc\'es du projet.

\subsection{Multicoeur}

\subsection{signaux}

\section{Conclusion}

\end{document}
